\begin{examplebox}
    \textbf{\textcolor{CExample}{例1（入门：绕原点旋转）}}
    \textbf{\textcolor{CExample}{题目：}} $B(0,0)$、$P(3,2)$。求 $P$ 分别绕 $B$ 顺时针和逆时针旋转 $90^\circ$ 得到的点，并验证垂直、等长。

    \textbf{\textcolor{CExample}{解题步骤：}}
    \begin{description}[style=nextline,leftmargin=2.8em,labelsep=0.5em,itemsep=0.4em]
        \item[\textbf{第一步（读出相对坐标）}]
          $\overrightarrow{BP}=(3,2)$，即从 $B$ 向右走 $3$、向上走 $2$。
        \item[\textbf{第二步（旋转）}]
          逆时针为 $(-2,3)$，顺时针为 $(2,-3)$；由于中心为原点，直接得到 $Q_1(-2,3)$、$Q_2(2,-3)$。
        \item[\textbf{第三步（验算两个条件）}]
          $\overrightarrow{BP}\cdot\overrightarrow{BQ_1}=3(-2)+2(3)=0$，$\overrightarrow{BP}\cdot\overrightarrow{BQ_2}=3(2)+2(-3)=0$。而 $BP^2=3^2+2^2=13$，$BQ_1^2=(-2)^2+3^2=13$，$BQ_2^2=2^2+(-3)^2=13$。
    \end{description}
    \textbf{\textcolor{CExample}{关键结论：}} 从原点旋转，直接交换坐标，并按转向改变符号。

    \medskip
    \textbf{\textcolor{CExample}{例2（平移中心：不能直接旋转绝对坐标）}}
    \textbf{\textcolor{CExample}{题目：}} 已知 $B(2,-1)$、$P(5,3)$，求满足 $BP\perp BQ$、$BP=BQ$ 的所有点 $Q$。

    \textbf{\textcolor{CExample}{解题步骤：}}
    \begin{description}[style=nextline,leftmargin=2.8em,labelsep=0.5em,itemsep=0.4em]
        \item[\textbf{第一步（先减旋转中心）}]
          $(u,v)=(5-2,3-(-1))=(3,4)$。
        \item[\textbf{第二步（两种旋转方向）}]
          逆时针相对坐标为 $(-4,3)$；顺时针相对坐标为 $(4,-3)$。
        \item[\textbf{第三步（加回旋转中心）}]
          逆时针 $Q_1=(2-4,-1+3)=(-2,2)$；顺时针 $Q_2=(2+4,-1-3)=(6,-4)$。
        \item[\textbf{第四步（检验距离）}]
          $BP^2=3^2+4^2=25$；$BQ_1^2=(-4)^2+3^2=25$，$BQ_2^2=4^2+(-3)^2=25$，且对应数量积均为 $0$。
    \end{description}
    \textbf{\textcolor{CExample}{关键结论：}} $\boxed{Q=(-2,2)\text{ 或 }Q=(6,-4)}$，不能把原点公式直接用在 $P(5,3)$ 上。

    \medskip
    \newpage
    \textbf{\textcolor{CExample}{例3（2020全国III卷：椭圆、定直线、面积定值）}}
    \textbf{\textcolor{CExample}{题目：}} 椭圆
    $\displaystyle C:\frac{x^2}{25}+\frac{y^2}{m^2}=1\ (0<m<5)$ 的离心率为 $\sqrt{15}/4$。左右顶点分别为 $A,B$。$P$ 在椭圆上，$Q$ 在 $x=6$ 上，且 $BP=BQ>0$、$BP\perp BQ$。求 $S_{\triangle APQ}$。

    \textbf{\textcolor{CExample}{解题步骤（优先旋转）}}
\begin{center}
\begin{tikzpicture}[x=0.65cm,y=0.65cm,>=stealth,every node/.style={font=\small},line cap=round]

\draw[->,gray] (-5.8,0)--(7.3,0) node[right] {$x$};
\draw[->,gray] (0,-1.9)--(0,3.2) node[above] {$y$};
\draw[thick] (0,0) ellipse (5 and 1.25);
\draw[gray,dashed] (6,-1.6)--(6,3.3) node[above] {$x=6$};
\fill[blue!8] (-5,0)--(3,1)--(6,2)--cycle;
\draw[blue!75!black,very thick] (5,0)--(3,1);
\draw[teal!75!black,very thick] (5,0)--(6,2);
\draw[thin] (-5,0)--(3,1)--(6,2)--cycle;
\draw (4.64,.18)--(4.82,.54)--(5.18,.36);
\draw[purple,dashed] (3,0)--(3,1) node[midway,left] {$1$};
\draw[purple,thick,<->] (5,-.65)--(6,-.65) node[midway,below] {$1$};
\fill (-5,0) circle (2pt);
\node at (-5,-1.65) {$A(-5,0)$};
\fill (5,0) circle (2pt);
\node at (3.8,-1.65) {$B(5,0)$};
\fill[blue!75!black] (3,1) circle (2pt);
\node[above left] at (3,2.1) {$P(3,1)$};
\fill[teal!75!black] (6,2) circle (2pt) node[above left] {$Q(6,2)$};
\node at (.5,-2.5) {画出一组解：横距 $6-5=1$ 对应 $P$ 的竖距 $|y_1|=1$。};

\end{tikzpicture}
\end{center}

    \begin{description}[style=nextline,leftmargin=2.8em,labelsep=0.5em,itemsep=0.4em]
        \item[\textbf{第一步（先求椭圆）}]
          由 $e^2=1-m^2/25=15/16$ 得 $m^2=25/16$，于是
          \[
             C:x^2+16y^2=25,\qquad A(-5,0),\ B(5,0).
          \]
        \item[\textbf{第二步（看见“垂直＋等长”，先交换坐标）}]
          设 $P(x_1,y_1)$、$Q(6,y_2)$。由于 $Q$ 相对 $B$ 的水平变化量是 $6-5=1$，而 $P$ 相对 $B$ 的竖直变化量为 $y_1$，立即有
          \[
             \boxed{|y_1|=|6-5|=1},\quad\text{即 } y_1=\pm1.
          \]
          再由椭圆方程 $x_1^2+16y_1^2=25$ 得 $x_1^2=9$，$x_1=\pm3$。
        \item[\textbf{第三步（同时确定 $Q$ 的纵坐标）}]
          若 $y_1=1$，只能顺时针旋转，$\overrightarrow{BQ}=(1,5-x_1)$，于是 $y_2=5-x_1$。若 $y_1=-1$，只能逆时针旋转，$\overrightarrow{BQ}=(1,x_1-5)$，于是 $y_2=x_1-5$。两种情况合写为
          \[
            \boxed{y_1y_2=5-x_1},\qquad\boxed{y_1^2=1}.
          \]
          因而无须把 $Q$ 的四个可能坐标一一代入面积公式。
        \item[\textbf{第四步（用坐标面积公式，保留整体结构）}]
          以 $A(-5,0)$ 为公共起点，可得
          \[
          \begin{aligned}
            S_{\triangle APQ}
            &=\frac12\left|(x_1+5)y_2-11y_1\right|\\
            &=\frac1{2|y_1|}\left|(x_1+5)(5-x_1)-11y_1^2\right|\\
            &=\frac1{2|y_1|}\left|25-x_1^2-11y_1^2\right|.
          \end{aligned}
          \]
          因为 $25-x_1^2=16y_1^2$，所以
          \[
            S_{\triangle APQ}=\frac1{2|y_1|}|5y_1^2|
             =\frac52|y_1|=\boxed{\frac52}.
          \]
    \end{description}
    \textbf{\textcolor{CExample}{回看原本的常规路线：}} 斜率条件得到 $y_2=(5-x_1)/y_1$；随后 $BQ^2=1+(5-x_1)^2/y_1^2=BP^2/y_1^2$，联立 $BP=BQ>0$ 才推出 $y_1^2=1$。旋转法直接越过了这段运算。

    \medskip
    \textbf{\textcolor{CExample}{例4（迁移：正方形相邻两边）}}
    \textbf{\textcolor{CExample}{题目：}} 正方形 $BPQR$ 按 $B\to P\to Q\to R$ 的次序排列，$B(1,2)$、$P(4,3)$。求其余顶点（注意有两种位置）。

    \textbf{\textcolor{CExample}{解题步骤：}}
\begin{center}
\begin{tikzpicture}[x=0.65cm,y=0.65cm,>=stealth,every node/.style={font=\small},line cap=round]

\fill[teal!8] (1,2)--(4,3)--(3,6)--(0,5)--cycle;
\draw[teal!75!black,thick] (1,2)--(4,3)--(3,6)--(0,5)--cycle;
\draw[orange!85!black,thick,dashed] (1,2)--(4,3)--(5,0)--(2,-1)--cycle;
\draw[blue!75!black,very thick,->] (1,2)--(4,3);
\draw[teal!75!black,very thick,->] (1,2)--(0,5);
\draw[orange!85!black,very thick,->] (1,2)--(2,-1);
\node[left] at (1,2) {$B(1,2)$}; \node[right] at (4,3) {$P(4,3)$};
\node[above] at (3,6) {$Q_1(3,6)$}; \node[left] at (0,5) {$R_1(0,5)$};
\node[right] at (5,0) {$Q_2(5,0)$}; \node[below] at (2,-1) {$R_2(2,-1)$};
\node[anchor=west] at (6,5) {旋转的是相邻边};
\node[anchor=west] at (6,3.8) {$BP\longrightarrow BR$};
\node[anchor=west] at (6,2) {由 $P$ 加上 $\overrightarrow{BR}$};
\node[anchor=west] at (6,.8) {才得到对角顶点 $Q$};

\end{tikzpicture}
\end{center}

    \begin{description}[style=nextline,leftmargin=2.8em,labelsep=0.5em,itemsep=0.4em]
        \item[\textbf{第一步（识别相邻边）}]
          在正方形中 $BP\perp BR$ 且 $BP=BR$，所以把 $\overrightarrow{BP}=(3,1)$ 绕 $B$ 转 $90^\circ$ 得到 $\overrightarrow{BR}$。
        \item[\textbf{第二步（求 $R$ 的两个位置）}]
          逆时针时 $\overrightarrow{BR}=(-1,3)$，$R_1=(0,5)$；顺时针时 $\overrightarrow{BR}=(1,-3)$，$R_2=(2,-1)$。
        \item[\textbf{第三步（利用平行四边形求 $Q$）}]
          正方形也是平行四边形，$\overrightarrow{PQ}=\overrightarrow{BR}$，于是 $Q_1=(4,3)+(-1,3)=(3,6)$；$Q_2=(4,3)+(1,-3)=(5,0)$。
    \end{description}
    \textbf{\textcolor{CExample}{关键结论：}} $\boxed{(Q,R)=((3,6),(0,5))\text{ 或 }((5,0),(2,-1))}$。适用的是\textbf{相邻两边} $BP,BR$，不是两条对角线。
\end{examplebox}
